\documentclass[../DoAn.tex]{subfiles}
\begin{document}
\begin{center}\bfseries\large TÓM TẮT NỘI DUNG BÀI TẬP LỚN\end{center}
\vspace{0.5cm}
RoomHub hỗ trợ giảng viên tìm và đăng ký phòng học, phòng họp; cán bộ quản trị xét duyệt và quản lý lịch sử dụng. Đề tài lấy quy trình ghi lịch trên bảng tính làm tình huống tham chiếu, từ đó xây dựng luồng tra cứu phòng, đặt phòng, phê duyệt, thông báo và cấp quyền tiếp cận. Dữ liệu thử nghiệm mô phỏng một tòa nhà tám tầng, lấy tòa C7 của Đại học Bách khoa Hà Nội làm bối cảnh tham chiếu; sơ đồ và danh mục phòng do nhóm xây dựng, không đại diện cho mặt bằng thực tế hoặc việc triển khai trên toàn trường.

Hệ thống gồm ứng dụng Android viết bằng React Native và TypeScript, dịch vụ Java 17/Spring Boot cung cấp REST API, cùng cơ sở dữ liệu H2 được truy cập qua JPA. Cấu hình hiện tại chọn chế độ dùng máy chủ: ứng dụng trao đổi dữ liệu qua API. Chế độ lưu cục bộ vẫn được duy trì riêng; dữ liệu ở hai chế độ không tự động hợp nhất khi thiết bị mất kết nối. Quy tắc nghiệp vụ được kiểm tra ở tầng dịch vụ; các thao tác ghi phía máy chủ chạy trong giao dịch có khóa bi quan.

Với phòng gắn khóa số, quản trị viên hoặc người dùng được cấp quyền theo phòng có thể tạo mã tạm thời gồm bảy chữ số. Khi có kết nối OneIoT trong phiên, ứng dụng gửi bản tin MQTT; kết quả gửi chưa chứng minh khóa thật đã áp dụng mã. Mã nguồn đã có dịch vụ máy chủ, cơ sở dữ liệu và luồng trao đổi với ứng dụng, nhưng khả năng chia sẻ dữ liệu giữa nhiều thiết bị trên hạ tầng triển khai cần được kiểm thử đầu cuối. Bản phát hành Android đã được kiểm tra chạy độc lập ở chế độ cục bộ.
\end{document}
